\documentclass[11pt,a4paper]{article}
\usepackage[T1]{fontenc}
\usepackage{lmodern,amsmath,amssymb,booktabs,microtype}
\usepackage[margin=24mm]{geometry}
\usepackage[colorlinks=true,urlcolor=blue,linkcolor=blue]{hyperref}
\newcommand{\Js}{\mathcal J_s}
\newcommand{\W}{\mathcal W_3}
\newcommand{\av}[1]{\langle #1\rangle}
\title{An independent cubic derivative singlet\\Gram projection and the first spectral moments}
\author{Alan Day \and Chengheng\thanks{AI research assistant (OpenAI).}}
\date{Local research note --- 6 November 2026}
\begin{document}\maketitle
\begin{abstract}
The cubic fermion bilinear contains an exact derivative of the quadratic
singlet. We isolate its Hermitian part, retain its third Gaussian chaos,
and project out the cubic energy polynomial and the energy-weighted
quadratic singlet. The remaining operator has strictly positive norm
and nonzero dynamics in complex DSSYK double scaling. Its second and
fourth frequency moments, and all nine two-arc moments through four
external Hamiltonians, agree with a probe of crossing weight $q^3$.
This gives a concrete candidate for the next independent singlet channel
and a finite-degree target for an all-order kernel reduction.
\end{abstract}
\section{Result and physical direction}
Use standard CAR $\{c_i,c_j^\dagger\}=\delta_{ij}$, a Hermitian
Hamiltonian $H=H_s$ with $p$ creations and $p$ annihilations in every
normal-ordered monomial, and the full annealed trace
$\av{Y}=\mathbb E_J 2^{-N}\operatorname{Tr}Y$ at
$\beta=\mu_{\rm chem}=0$. The complex Gaussian ensemble includes
overlapping index sets. First take $N,p\to\infty$ at fixed
\[
 \lambda=p^2/N>0,\qquad q=e^{-4\lambda},\qquad
 \sigma^2=\av{H^2}=\Js^2/\lambda .
\]
Write $L=[H,\,\cdot\,]$, $A=\sum_i(Lc_i)^\dagger Lc_i$,
$M_3=\sum_i c_i^\dagger L^3c_i$, and
$T_3=(M_3+M_3^\dagger)/2$. Our operator is
\begin{equation}
 R=T-aS-bZ,\quad
 (T,S,Z)=\W\bigl(T_3,\{H,A\}/2,H^3\bigr),                 \label{R}
\end{equation}
where $\W$ means third \emph{Gaussian coupling chaos}, not fermion
normal ordering. The coefficients are
\begin{equation}
 a=-\frac{(1-q)^2}{1+q^2},\qquad
 b=\frac{2\lambda(1-q)(1+q^3)}
              {(1+q^2)(1+q+q^2)}.                       \label{ab}
\end{equation}
For $\widehat R=R/\sqrt{\av{R^2}}$, the frequency moments are
\begin{align}
 \mu_2&=\av{\widehat R L^2\widehat R}
       =\frac{2\Js^2}{\lambda}(1-q^3),                 \label{mu2}\\
 \mu_4&=\frac{2\Js^4}{\lambda^2}(1-q)^2(1+q+q^2)
                   (3q^3+6q^2+6q+5).                 \label{mu4}
\end{align}
Thus $\mu_2\to24\Js^2$ and $\mu_4\to1920\Js^4$ as
$\lambda\to0^+$. These are microscopic frequency data for the
Susskind singlet program \cite{MSS}; they do not yet fix a bulk mass.

\newpage
\section{Removing the exact time descendant}
Define $B_{ab}=\sum_i(L^ac_i)^\dagger L^bc_i$. Since
$(LY)^\dagger=-L(Y^\dagger)$, the product rule gives
\begin{equation}
 LB_{ab}=-B_{a+1,b}+B_{a,b+1}.                           \label{rec}
\end{equation}
Homogeneity in the fermions implies $B_{01}=-pH$, so
$B_{02}=B_{11}=A$. Applying (\ref{rec}) twice gives
\begin{align}
 M_3&=2B_{12}-B_{21},&
 M_3^\dagger&=2B_{21}-B_{12},\\
 T_3&=\tfrac12(B_{12}+B_{21}),&
 \boxed{M_3=T_3+\tfrac32 LA.}                         \label{split}
\end{align}
These are exact finite-$N$ identities for each realization. In real
time $\partial_t^n c=i^nL^nc$; retaining these factors is essential
when comparing the schematic derivative family in \cite{MSS}.

For a direct construction let $Q_f=\sum_i c_i^\dagger c_i$ and
\[
 \mathfrak D(Y)=Q_fY-\sum_i c_i^\dagger Yc_i .
\]
On a normal-ordered number-conserving monomial with $k$ creation
operators, $\mathfrak D$ multiplies by $k$. Expanding the commutators
therefore gives
\begin{align}
 A&=(2p-\mathfrak D)H^2,\\
 T_3&=(3p-\mathfrak D)H^3-\tfrac32\{H,A\}.             \label{rank}
\end{align}
This also explains the contraction marks used below. For three
individually $p$-body factors $V_1,V_2,V_3$, set
\begin{align}
 K_{12}&=((2p-\mathfrak D)(V_1V_2))V_3,\\
 K_{23}&=V_1((2p-\mathfrak D)(V_2V_3)),\\
 K_{13}&=(3p-\mathfrak D)(V_1V_2V_3)-K_{12}-K_{23}.
\end{align}
Each normal-order contraction reduces the total body rank by one.
Consequently $K_{ij}$ counts contractions between those two factors,
including the nonadjacent pair $13$. With all $V_i=H$,
\begin{equation}
 T_3=K_{13}-\tfrac12(K_{12}+K_{23}),\qquad
 \{H,A\}/2=\tfrac12(K_{12}+K_{23}).                    \label{marks}
\end{equation}

There is also an exact ensemble selection rule:
$\av{T_3L^nA}=0$ for every $n$. For even $n$, use $H\mapsto-H$.
For odd $n$, the trace is purely imaginary by Hermiticity, but its
average is real by complex conjugation of the couplings in the real
occupation basis. The same argument applies to $R$ and the previously
constructed even quadratic singlet $X$. This removes the simple
$X$ time-descendant sector; energy-weighted composites still require
the explicit projection (\ref{R}).

\newpage
\section{A positive residual from the six-slot Gram form}
For Gaussian coordinates $J_\alpha$ with covariance $C_{\alpha\beta}$,
\[
 \W(J_\alpha J_\beta J_\gamma)
 =J_\alpha J_\beta J_\gamma
 -C_{\alpha\beta}J_\gamma-C_{\alpha\gamma}J_\beta
 -C_{\beta\gamma}J_\alpha .
\]
When computing a product of such insertions, Wick pairings internal
to either three-slot insertion are excluded. This specifies the
operator before taking double scaling. The lower, linear chaos of
the raw $T_3$ has not been identified here with a conserved operator.

The Gram matrix $G_{ij}=\av{Y_iY_j}/\sigma^6$, for $Y=(T,S,Z)$,
is the sum of the six cross pairings. Its independent entries are
\begin{align}
 G_{TT}&=2\lambda(4\lambda q^2+4\lambda q+q^2-4q+3),\\
 G_{TS}&=-2\lambda(q-1)^2,&G_{TZ}&=-4\lambda q(q-1),\\
 G_{SS}&=2\lambda(2\lambda q^3+2\lambda+q^2+1),\\
 G_{SZ}&=-2\lambda(q-1)(q^2+q+1),\\
 G_{ZZ}&=(q+1)(q^2+q+1).
\end{align}
Solving $\av{RS}=\av{RZ}=0$ yields (\ref{ab}). Define
\[
 \mathcal B=1-q^3+4\lambda q^2(2q^2-q+2).
\]
The Schur complement is
\begin{equation}
 V:=\frac{\av{R^2}}{\sigma^6}
   =\frac{4\lambda(1+q)\mathcal B}{(1+q^2)(1+q+q^2)}
   =32\lambda^2-96\lambda^3+O(\lambda^4).              \label{V}
\end{equation}
Since $0<q<1$ and $2q^2-q+2>0$, $V>0$. In physical units
$\av{R^2}\sim32\Js^6/\lambda$. The direction therefore survives
normalization even though its dimensionless Gram weight is small.

\paragraph{How the sums are evaluated.}
For a pairing involution $\pi$ on $\ell$ Hamiltonian slots and two
marks $m_1,m_2$, form the local entries
\[
 E_{uv}=\operatorname{sgn}(\pi(v)-u+\tfrac12)
       -2x\,\mathbf1_{(\pi(v),u)=m_1}
       -2y\,\mathbf1_{(\pi(v),u)=m_2}.
\]
Expand principal two-by-two minors with $x^2=y^2=0$ and divide by
their singleton signs. Adding the normalization contribution
$-\binom{\ell}{2}-\ell/2$ gives
$-4\,{\rm cr}+\alpha x+\beta y+\delta xy$.
The fixed-degree double-scaling weight is consequently
\[
 q^{\rm cr}(\lambda\delta+\lambda^2\alpha\beta)
\]
for two marks, $q^{\rm cr}\lambda\alpha$ for one, and $q^{\rm cr}$
for none. This follows from the Poisson pair-overlap limit; higher
mode coincidences vanish at fixed degree. Equations (\ref{marks})
then give the displayed matrix. The nonadjacent mark is checked
separately against finite CAR products, rather than inferred from
an independent-probe formula.

\newpage
\section{Dynamics and a finite-degree kernel reduction}
Define $K_{ij}=\av{(LY_i)^\dagger LY_j}/\sigma^8$.
Cyclic trace gives the computable real symmetric expression
\begin{equation}
 \sigma^8K_{ij}
 =\av{Y_iY_jH^2}+\av{Y_jY_iH^2}-2\av{Y_iHY_jH}.
\end{equation}
All Gaussian pairings are retained except those internal to a
Wick-ordered insertion. There are 42 allowed pairings in each
eight-slot arrangement, not just $4!$ fully cross pairings.
The nonzero entries are
\begin{align}
 K_{TT}&=2\lambda(1-q^2)
 [4\lambda q^3+16\lambda q^2+4\lambda q+q^3-3q^2-3q+5],\\
 K_{TS}&=2\lambda(q-1)^3(q+1)(4\lambda q+q+1),\\
 K_{SS}&=2\lambda(1-q^2)(1+q^2)(4\lambda q+q+1).
\end{align}
The $Z$ row vanishes. For $v=(1,-a,-b)^T$ one obtains the
particularly simple exact identity
\begin{equation}
 v^TKv=2(1-q^3)V,\qquad
 v^TKv=768\lambda^3+O(\lambda^4).                      \label{ratio}
\end{equation}
This proves (\ref{mu2}) and nonzero dynamics at every $\lambda>0$.

More information is available than this single ratio. Let
\[
 F_{uv}=\frac{\av{RH^uRH^v}}{\av{R^2}\sigma^{u+v}} .
\]
Exact marked enumeration gives the following complete set through
four external Hamiltonians (with $F_{uv}=F_{vu}$):
\begin{center}
\begin{tabular}{ccl}
\toprule
$u+v$ & $(u,v)$ & $F_{uv}$\\
\midrule
0 & $(0,0)$ & $1$\\
2 & $(2,0)$ & $1$\\
2 & $(1,1)$ & $q^3$\\
4 & $(4,0)$ & $2+q$\\
4 & $(3,1)$ & $q^3(2+q)$\\
4 & $(2,2)$ & $1+q^6+q^7$\\
\bottomrule
\end{tabular}
\end{center}
All corresponding $\av{RH^uSH^v}$ and
$\av{RH^uZH^v}$ vanish. Combining the fourth-order row with
the binomial expansion of $L^4$ gives (\ref{mu4}).

These nine two-arc values coincide with the existing independent
probe rule \cite{BNR}: each Hamiltonian crossing of the probe carries
$r=q^3$. This is a tested finite-degree reduction for a
\emph{coupling-dependent composite}; independence of its couplings
has not been assumed. An all-order proof would permit use of the
known $q$-Mehler two-point kernel. It would also predict
$\av{\widehat R(t)\widehat R(0)}\to\operatorname{sech}^{6}(2\Js t)$
at fixed string time as $\lambda\to0^+$. At this stage that full
function is a conjecture supported by the computed moments.

\newpage
\section{Evidence, provenance, and the next calculation}
The motivation is the derivative singlet family in
Miyashita--Sekino--Susskind, Sec.~4.3 \cite{MSS}. Their proposed
$n=3$ assignment motivates this test but does not supply the
quantitative correlator. Berkooz--Narovlansky--Raj,
Sec.~7 and Eqs.~(7.1)--(7.3) \cite{BNR}, solve a probe whose random
couplings are uncorrelated with those of $H$. We use their kernel
as a comparison target; the new work here is the explicit
same-coupling composite and its finite-degree reduction.
No priority claim is made for the polynomial identities.

The saved checks comprise:
\begin{itemize}
\item Twelve operator identities in independent exact
integer-complex SymPy CAR matrices for $(N,p)=(3,1),(4,2)$.
The same identities also pass the direct occupation-basis code.
\item 192 fixed six-slot words, including coherent repeated index
sets, with all 16 choices of zero, adjacent, or nonadjacent marks:
3072 finite CAR/determinant comparisons.
\item Six exact symbolic projection, positivity-form, kinetic,
and small-$\lambda$ certificates; nine two-arc equalities and
18 vanishing cross moments through degree four.
\end{itemize}
The analytic positivity follows from (\ref{V}), rather than from
a parameter scan. Double scaling is justified at fixed word degree;
finite CAR checks test the contraction interpretation and signs,
not convergence in $N$. The scripts record dependencies and input
hashes. A parsing failure caused by the reserved word
\texttt{lambda} was fixed before the higher-moment output was
written; no earlier scientific output was replaced.

\paragraph{Next research step.}
Prove or disprove the $q^3$ reduction for arbitrary two-arc words.
The useful target is cancellation of diagrams in which one of the
three insertion slots contracts with an external $H$ after the
$S,Z$ subtraction. A finite local endpoint identity, analogous to
the earlier quadratic construction, could reduce the remaining
three lines to one probe line. This is preferable to merely
extending the moment table. In parallel with that analytic step,
the linear Gaussian chaos of raw $T_3$ should be identified so
that the relation to the unprojected derivative operator is complete.

The present positive residual establishes independence from the
displayed cubic composites and from the simple $X$ time-descendant
sector in the ensemble inner product. A complete particle
classification still requires the full operator dictionary and
spatial dynamics. In particular, neither the derivative index nor
the $r=q^3$ candidate label is by itself a mass eigenvalue.
Local manuscript and PDF only; no Git commit, push, or archive
was made in this stage.

\begin{thebibliography}{9}
\bibitem{MSS}
S. Miyashita, Y. Sekino and L. Susskind,
\emph{Holograms and Standard Models},
\href{https://arxiv.org/html/2607.05678v2}{arXiv:2607.05678v2},
Sec.~4.3, Eqs.~(4.19)--(4.26).
\bibitem{BNR}
M. Berkooz, V. Narovlansky and H. Raj,
\emph{Complex Sachdev--Ye--Kitaev model in the double scaling limit},
\href{https://arxiv.org/html/2006.13983v2}{arXiv:2006.13983v2},
Sec.~7, Eqs.~(7.1)--(7.3).
\end{thebibliography}
\end{document}
